% ED26-SRC-063-0104: original Johnson 2004 continuation, pp.265--307.
% Follows the frozen Johnson foundation; no second chapter heading.
\section{McCall's native categorical modal calculus}
\label{sec:p13-johnson-mccall}

Johnson's third approach is not a new choice of S5 predicate translation.
McCall retains categorical expressions and propositional assertion, adds
restricted modality, and separately axiomatizes rejection. Johnson subsequently
supplies a term-region semantics for the syntactic calculus. The syntactic
McCall system, Johnson's semantics, and the later contingency modification
must not be identified with each other
\parencite[pp.~265--276]{Johnson2004ModalSyllogisms}.

\begin{hbdefinition}{def:p13-johnson-mccall-grammar}{Categorical modal formation}
Term variables are \(a,b,c,\ldots\). A categorical expression is an \(Aab\)
or \(Iab\), or a negation of a categorical expression. Categorical expressions
are formulas, and \(Lp\) is a formula only when \(p\) is a categorical
expression. Formulas are closed under \(\neg\) and \(\to\), with no other
formation cases. Keep the subject-first definitions
\(Eab\mathrel{:\equiv}\neg Iab\),
\(Oab\mathrel{:\equiv}\neg Aab\), and
\(Mp\mathrel{:\equiv}\neg L\neg p\).
Thus \(L\neg\neg Aab\) is permitted, whereas \(LLAab\) and
\(L(Aab\to Iab)\) are not formulas of this language.
\end{hbdefinition}

\(X\) in a mood code is metadata for an unmodalized categorical formula,
not an object-language operator. A code such as LXL records the major premise,
minor premise, and conclusion modes in that order. \(\vdash p\) denotes
assertion in the indicated system; \(\dashv p\) denotes rejection. Rejection
does not assert \(\neg p\). The existential axiom in L-X-M is
\(\vdash LIaa\), stronger than the earlier assertoric \(\vdash Iaa\).
The new contingency calculus will deliberately restore \(Iaa\) instead.

\subsection{Assertion basis and legitimate transformations}
\label{sec:p13-johnson-mccall-basis}

The assertion basis retains propositional tautologies A0, \(Aaa\) (A1),
assertoric Barbara \(Abc\to(Aab\to Aac)\) (A3), and assertoric Datisi
\(Abc\to(Iba\to Iac)\) (A4). A2 is now \(LIaa\).
Its additional axioms are the following original subject-first formulas.

\begin{tabularx}{\textwidth}{@{}l l X@{}}
\toprule
Axiom & Interpretation & Formula asserted\\
\midrule
A5 & Barbara LXL & \(LAbc\to(Aab\to LAac)\)\\
A6 & Cesare LXL & \(LEcb\to(Aab\to LEac)\)\\
A7 & Darii LXL & \(LAbc\to(Iab\to LIac)\)\\
A8 & Ferio LXL & \(LEbc\to(Iab\to LOac)\)\\
A9 & Baroco LLL & \(LAcb\to(LOab\to LOac)\)\\
A10 & Bocardo LLL & \(LObc\to(LAba\to LOac)\)\\
A11 & Necessary particular conversion & \(LIab\to LIba\)\\
A12 & Universal subordination & \(LAab\to Aab\)\\
A13 & Particular subordination & \(LIab\to Iab\)\\
A14 & Negative particular subordination & \(LOab\to Oab\)\\
\bottomrule
\end{tabularx}

Uniform term substitution, modus ponens, and definition interchange remain.
AR4 additionally permits inserting or removing two negations on a categorical
expression within a well-formed assertion. This is not a rule to necessitate
an arbitrary asserted formula. Classical propositional reasoning treats the
well-formed modal categoricals as propositional units; it does not move a
propositional equivalence through \(L\) indiscriminately. The restricted
double-negation operation is what licenses the required modal oppositions
\parencite[pp.~265--268]{Johnson2004ModalSyllogisms}.
The literal operations are in \cref{sheet:p13-johnson-mccall}.
\calculussheet{sheet:p13-johnson-mccall}{Johnson's native McCall operations}{Subject-first modal categoricals; separate assertion and rejection; L and Q take only categorical expressions; no nested modality or individual equality calculus.}
\hbrule{rule:p13-johnson-mccall-us}{AR1 Uniform term substitution}{\vdash p}{\vdash p\sigma}{sigma substitutes term variables uniformly; resulting expression well formed.}
\hbrule{rule:p13-johnson-mccall-mp}{AR2 Modus ponens}{\vdash(p\to q)\quad\vdash p}{\vdash q}{Same indicated assertion system.}
\hbrule{rule:p13-johnson-mccall-def}{AR3 Definition interchange}{\vdash C[\alpha]}{\vdash C[\beta]}{alpha and beta are a definiendum and its declared definiens; either direction; C remains well formed.}
\hbrule{rule:p13-johnson-mccall-dn}{AR4 Categorical double negation}{\vdash C[p]}{\vdash C[\neg\neg p]}{p is a categorical expression; either direction; C remains well formed.}
\hbrule{rule:p13-johnson-mccall-rus}{RR1 Rejection substitution}{\dashv p\sigma}{\dashv p}{sigma is a uniform term substitution; the rejected instance and general formula are well formed.}
\hbrule{rule:p13-johnson-mccall-rd}{RR2 Rejection detachment}{\vdash(p\to q)\quad\dashv q}{\dashv p}{Same indicated assertion/rejection system; not assertion of neg p.}
\hbrule{rule:p13-johnson-mccall-rs}{RR3 Restricted Slupecki rule}{\dashv(x\to z)\quad\dashv(y\to z)}{\dashv(x\to(y\to z))}{x and y are neg Auv or neg Iuv; z is a nested elementary conditional with only simple A I E O categorical components.}
\hbrule{rule:p13-johnson-mccall-rdef}{RR4 Rejection definition interchange}{\dashv C[\alpha]}{\dashv C[\beta]}{Declared definiendum/definiens; either direction; C remains well formed.}
\hbrule{rule:p13-johnson-mccall-rdn}{RR5 Rejection categorical double negation}{\dashv C[p]}{\dashv C[\neg\neg p]}{p is categorical; either direction; C remains well formed.}
\hbrule{rule:p13-johnson-mccall-rweak}{Derived rejection antecedent weakening}{\dashv(p\to(q\to r))\quad\vdash(q\to s)}{\dashv(p\to(s\to r))}{Derived from A0 and RR2; similarly for the first antecedent.}
\hbrule{rule:p13-johnson-mccall-rstrong}{Derived rejection consequent strengthening}{\dashv(p\to(q\to r))\quad\vdash(s\to r)}{\dashv(p\to(q\to s))}{Derived from A0 and RR2.}


\begin{hbtheorem}{thm:p13-johnson-mccall-immediate}{proved=annex:p13-johnson-mccall-immediate}{Native modal immediate inferences}
The basis just specified asserts the four opposition equivalences
\(\neg LAab\leftrightarrow MOab\),
\(\neg LEab\leftrightarrow MIab\),
\(\neg LIab\leftrightarrow MEab\), and
\(\neg LOab\leftrightarrow MAab\), and their negated reverse forms.
It asserts conversion for \(LE,LI,ME,MI\); subalternation
\(LA\to LI\), \(LE\to LO\), \(MA\to MI\), \(ME\to MO\);
and conversion per accidens \(LAab\to LIba\), \(LEab\to LOba\),
\(MAab\to MIba\), \(MEab\to MOba\).
For each categorical kind \(T\in\{A,E,I,O\}\), it also asserts
\(LTab\to Tab\) and \(Tab\to MTab\).
\end{hbtheorem}

The theorem's complete native derivations are in the annex. In particular,
Johnson's printed theorem~4.3 first line cites A11 for
\(LAab\to LIba\), although A11 only gives \(LIab\to LIba\).
The annex derivation uses A7 with \(Iaa\), obtained from A2 and A13,
then A11. This records and repairs the missing dependency rather than
admitting the desired subalternation as a new axiom.

The distinction between Barbara LXL and XLL is now substantive. A5 asserts
the former, but XLL is rejected below. Classical contraposition of a nested
conditional generates other mode combinations without an unrestricted modal
distribution law. For instance, reversing a Ferison LXL instance yields Darii
MXM. Necessity on the first categorical premise therefore cannot be moved
to the second merely because the unmodalized syllogism is valid.

\section{Rejection in the native system}
\label{sec:p13-johnson-mccall-rejection}

Johnson explicitly gives only part of McCall's rejection basis. R1 remains
the rejected assertoric formula \(Acb\to(Aab\to Iac)\).
The three further rejection axioms printed on p.~266 are:
\[
\begin{aligned}
\sigma_2&=LAbb\to(MAab\to(Aac\to(LAca\to(LAbc\to LAac)))),\\
\sigma_3&=LAaa\to(LAcc\to(MAac\to(LAca\to Aac))),\\
\sigma_4&=LAaa\to(LAbb\to(LAcc\to(LAab\to
 (MAba\to(MAbc\to(LAcb\to Iac)))))).
\end{aligned}
\]
R2--R4 reject these respective formulas. In \(\sigma_3\), the final
consequent is \(Aac\), not \(LAac\); Johnson's note~15 expressly corrects
the earlier 1989 display. These are rejection axioms, not false assertions
that their negations are universally true.

Rejection substitution runs backward: rejection of a uniform term-substitution
instance rejects the general formula. Rejection detachment combines an asserted
\(p\to q\) with rejected \(q\) to reject \(p\). Definition interchange and
categorical double-negation interchange also apply. S\l{}upecki's RR3 is
restricted: its two antecedents must have the forms \(\neg Axy\) or
\(\neg Ixy\), and its common consequent must be an elementary nested
conditional of simple categoricals. It is not an unrestricted rule for
combining modal rejections.

Classical equivalent conditionals permit rejection reversal and antecedent
interchange through rejection detachment. Rejection antecedent weakening and
consequent strengthening point in the opposite directions from the assertion
operations. All such steps used in the following result are justified locally
in the annex, not by a bare reference to a historical theorem name
\parencite[pp.~253--254, 266, 270--271]{Johnson2004ModalSyllogisms}.

\begin{hbtheorem}{thm:p13-johnson-mccall-rejections}{proved=annex:p13-johnson-mccall-rejections}{Ten separating syntactic rejections}
Using the displayed R2--R4 and the stated transformations, L-X-M rejects
Barbara XLL; Baroco XMM and Bocardo MLX; Barbara LMX, Baroco LXL,
and Bocardo XMM; Barbari MLX, Bramantip LMX, Felapton XLL, and Bocardo XLL.
\end{hbtheorem}

The source's theorem~4.12 last formula is
\(Obc\to(LAba\to LOac)\), which is \emph{Bocardo XLL}, not Baroco XLL.
The result above therefore uses that actual formula. Baroco XLL has a separate
semantic countermodel below; its syntactic rejection is not invented from
the partial rejection basis. Names, modes, and literal formulas must be
checked together. The source's printed labels and citations are not
additional inference rules.

\section{Johnson's two-layer semantics for L-X-M}
\label{sec:p13-johnson-mccall-semantics}

\begin{hbdefinition}{def:p13-johnson-jmodel}{The native term-region model}
A J-model has a nonempty object domain \(W\). For each term \(x\), the four
sets \(n^{+}(x),q^{+}(x),n^{-}(x),q^{-}(x)\) partition \(W\).
They represent essential positive, contingent positive, essential negative,
and contingent negative regions. Put
\(+(x)=n^{+}(x)\cup q^{+}(x)\), and write \(U\mathbin{o}V\) for
\(U\cap V\ne\varnothing\). The source's additional base conditions are:
\[
\begin{array}{ll}
\mathrm{B2}&n^{+}(x)\ne\varnothing,\\
\mathrm{B3}&+(z)\subseteq n^{-}(y),\ +(x)\subseteq +(y)
              \ \Longrightarrow\ +(x)\subseteq n^{-}(z),\\
\mathrm{B4}&+(y)\subseteq n^{+}(z),\ +(x)\mathbin{o}+(y)
              \ \Longrightarrow\ n^{+}(x)\mathbin{o}n^{+}(z),\\
\mathrm{B5}&+(y)\subseteq n^{-}(z),\ +(x)\mathbin{o}+(y)
              \ \Longrightarrow\ n^{+}(x)\mathbin{o}n^{-}(z),\\
\mathrm{B6}&+(z)\subseteq n^{+}(y),\ n^{+}(x)\mathbin{o}n^{-}(y)
              \ \Longrightarrow\ n^{+}(x)\mathbin{o}n^{-}(z).
\end{array}
\]
All term parameters are universal in these base conditions. B2 requires an
essential positive instance for \emph{every term}, not merely a nonempty universe.
\end{hbdefinition}

These are not sets of accessible worlds. They are four regions of one object
domain for each term. The superstructural truth conditions are two-valued:
\[
\begin{array}{lll}
Axy:&+(x)\subseteq +(y),& Ixy:\ +(x)\mathbin{o}+(y),\\
LAxy:&+(x)\subseteq n^{+}(y),& LIxy:\ n^{+}(x)\mathbin{o}n^{+}(y),\\
LOxy:&n^{+}(x)\mathbin{o}n^{-}(y),& LExy:\ +(x)\subseteq n^{-}(y).
\end{array}
\]
Negation and implication use ordinary Boolean truth, and categorical
double-negation under \(L\) leaves truth unchanged. Thus \(M\) is evaluated
by its definition, not by a new accessible-world relation
\parencite[pp.~271--276]{Johnson2004ModalSyllogisms}.

The printed definition~5.2 says \(V(p\to q)=t\) when \(V(\neg p)=t\)
or \(V(q)=f\). The latter value must be \(t\), as required by the chapter's
classical tautologies, detachment, and countermodels. The active semantics
uses the corrected consequent value. This is a recorded formula correction,
not an alternative implication silently attributed to McCall.

\subsection{Four original countermodels and their separating cases}
\label{sec:p13-johnson-mccall-countermodels}

The following are the four source models from tables~6,9,10,11. Braces and
empty sets are explicit here. Unmentioned term variables receive the \(a\)
row, as in Johnson's convention. The domain of each is the union of one
row's four regions, not a further set of possible worlds.

\begin{longtable}{@{}llcccc@{}}
\caption{Johnson's four source countermodels (source tables~6, 9, 10 and 11) with every region explicit}\\
\toprule Model & Term & \(n^{+}\) & \(q^{+}\) & \(n^{-}\) & \(q^{-}\)\\
\midrule\endfirsthead
\caption[]{Johnson's four source countermodels (source tables~6, 9, 10 and 11) with every region explicit (continued)}\\
\toprule Model & Term & \(n^{+}\) & \(q^{+}\) & \(n^{-}\) & \(q^{-}\)\\
\midrule\endhead
\midrule\multicolumn{6}{r}{\emph{continued on next page}}\\
\endfoot
\endlastfoot
\(\mathcal M_1\)&a&\(\{1\}\)&\(\{2\}\)&\(\varnothing\)&\(\{3\}\)\\
&b&\(\{1,2,3\}\)&\(\varnothing\)&\(\varnothing\)&\(\varnothing\)\\
&c&\(\{3\}\)&\(\{1,2\}\)&\(\varnothing\)&\(\varnothing\)\\
\(\mathcal M_2\)&a&\(\{1\}\)&\(\varnothing\)&\(\{2\}\)&\(\{3\}\)\\
&b&\(\{1\}\)&\(\varnothing\)&\(\{2\}\)&\(\{3\}\)\\
&c&\(\{2\}\)&\(\{3\}\)&\(\varnothing\)&\(\{1\}\)\\
\(\mathcal M_3\)&a&\(\{1\}\)&\(\{2\}\)&\(\varnothing\)&\(\{3\}\)\\
&b&\(\{3\}\)&\(\varnothing\)&\(\{2\}\)&\(\{1\}\)\\
&c&\(\{3\}\)&\(\varnothing\)&\(\{2\}\)&\(\{1\}\)\\
\(\mathcal M_4\)&a&\(\{1,2\}\)&\(\{3,4\}\)&\(\varnothing\)&\(\varnothing\)\\
&b&\(\{1,3\}\)&\(\{4\}\)&\(\{2\}\)&\(\varnothing\)\\
&c&\(\{1\}\)&\(\{4\}\)&\(\{3\}\)&\(\{2\}\)\\
\bottomrule
\end{longtable}

\begin{hbtheorem}{thm:p13-johnson-mccall-models}{proved=annex:p13-johnson-mccall-models}{Native model validation and six exact separations}
These four structures satisfy B1--B6. Respectively they invalidate Barbara
XLL, Celarent XLL, Camestres LXL, and Baroco XLL. Row permutations preserve
the model conditions; permutations of \(\mathcal M_1\) also invalidate
Ferio MLX and Baroco XMM. Each case has both categorical antecedents true
and its conclusion false.
\end{hbtheorem}

Johnson's table~5 uses codes such as \(1ac\) to specify one of these models
and a placement of its rows into minor and major term roles. A blank is an
asserted combination in his comparison, not an unassessed entry. The source's
initial sentence calls \(\mathcal M_1\) a countermodel to Barbara LXL,
but its computation and cell are Barbara \emph{XLL}; A5 asserts LXL.
Likewise the p.~274 Ferio calculation contains a misindexed \(MEac\)
and \(Obc\): the genuine subject-first case is
\(MEbc,LIab;Oac\). The annex computes the actual formulas, rather than
certifying those slips from the table's mood name
\parencite[pp.~273--277]{Johnson2004ModalSyllogisms}.

The four countermodels and their permutations do not by themselves prove
general completeness or provide countermodels for every rejection axiom.
Johnson expressly gives an additional model for McCall's more complicated
rejection *5.41 and reports a six-object bound for certain finite nested
conditional forms from his 1989 result. The full missing rejection basis and
that general bound remain cited source dependencies, not conclusions of six
examples.

\subsection{Thomason's stronger conditions and the chain boundary}
\label{sec:p13-johnson-mccall-t3-chain}

Thomason's T3 presentation keeps the partition and essential-positive
nonemptiness, but replaces B3--B6 by three conditions:
\[
\begin{array}{ll}
\mathrm{BT3}&+(x)\mathbin{o}+(y)\ \Longrightarrow\ +(x)\mathbin{o}n^{+}(y),\\
\mathrm{BT4}&+(x)\subseteq n^{-}(y)\ \Longrightarrow\ +(y)\subseteq n^{-}(x),\\
\mathrm{BT5}&+(x)\subseteq n^{+}(y)\ \Longrightarrow\ n^{-}(y)\subseteq n^{-}(x).
\end{array}
\]

\begin{hbtheorem}{thm:p13-johnson-t3-comparison}{proved=annex:p13-johnson-t3-comparison}{The exact model-class comparison}
Every T3-model is a J-model. The converse fails: the displayed
\(\mathcal M_1\) is a J-model but not a T3-model. B3 and BT4 are
equivalent when universally quantified over their term parameters.
\end{hbtheorem}

The source's theorem~5.6 uses \(\mathcal M_4\) for the strictness example,
but its rows all have 1 in both the actual positive and essential positive
regions, so its claimed BT3 failure does not occur. The comparison above uses
\(\mathcal M_1\): actual a overlaps actual c, but not essential c.
This correction affects the example, not the valid inclusion argument
\parencite[pp.~278--279]{Johnson2004ModalSyllogisms}.

Johnson then distinguishes nonaxiomatic variants with singular sentences and
ecthesis from McCall's axiomatic L-X-M. Darapti's setting-out selects an
existent b, propagates its actual membership through Abc and Aba, and
generalizes the shared a/c witness. The apodeictic Baroco setting-out needs
an additional modal witness condition; the necessary nonmembership propagation
is not ordinary existential instantiation alone. The source's p.~280 line 5
explanation mistakenly says necessarily not c where its input is necessarily
not b. These variants and their completeness results remain separately cited,
not adopted through that slipped explanatory sentence.

For the relevance-oriented AP system, the premises together with the
categorical contradictory of the proposed conclusion form a simple cycle of
term links. Distinct vertices are required; \(\{LAaa\}\) and repeated
edges such as \(\{LAab,Aba,MAac,Aca\}\) are excluded. Ignoring direction,
the chain links the conclusion terms; it is not an arbitrary set of classical
antecedents. The contradictory map pairs A/O and I/E, and pairs LA/MO,
LE/MI, LI/ME, and LO/MA, with the reverse map involutive.
Johnson gives eleven accepted cycle forms for AP and three assertoric
restrictions, under the validity definitions of his cited 1994 source and
Smiley's cited 1973 source. For example, the cycle
\(\{LAab,LAbc,Acd,Oad\}\) certifies the corresponding Aad consequence
in that reported procedure. These decision procedures are separate from
the scoped derivations above; their relevance condition is not a condition
of the unrestricted classical assertion basis
\parencite[pp.~279--284]{Johnson2004ModalSyllogisms}.

\section{Contingency: a new primitive and a failed extension}
\label{sec:p13-johnson-mccall-contingency}

Johnson distinguishes one-sided possibility, not impossible, from two-sided
contingency, neither necessary nor impossible. Prior's familiar explanation
uses \(Mp\land M\neg p\), but McCall's \(Q\) is a \emph{primitive}
categorical operator. The system does not simply define every \(Qp\) by that
modal-propositional conjunction. In particular, complementary conversion
changes categorical quality, whereas ordinary conversion exchanges terms.
Conflating the two would validate the disputed universal contingent converse
\parencite[pp.~283--286]{Johnson2004ModalSyllogisms}.

The Q-L-X-M formation rule permits \(Qp\) only for categorical expressions,
just as for \(Lp\). It uses the earlier assertion operations, changes A2
\emph{back} to \(Iaa\), and adds these original axioms:

\begin{longtable}{@{}llp{0.59\textwidth}@{}}
\caption{Contingency: a new primitive and a failed extension}\\
\toprule Axiom & Name & Formula\\\midrule\endfirsthead
\caption[]{Contingency: a new primitive and a failed extension (continued)}\\
\toprule Axiom & Name & Formula\\\midrule\endhead
\midrule\multicolumn{3}{r}{\emph{continued on next page}}\\
\endfoot
\endlastfoot
A15 & Barbara QQQ & \(QAbc\to(QAab\to QAac)\)\\
A16 & Darii QQQ & \(QAbc\to(QIab\to QIac)\)\\
A17 & First-figure QXQ-AAE & \(QAbc\to(Aab\to QEac)\)\\
A18 & Darii QXQ & \(QAbc\to(Iab\to QIac)\)\\
A19 & Barbara XQM & \(Abc\to(QAab\to MAac)\)\\
A20 & Celarent XQM & \(Ebc\to(QAab\to MEac)\)\\
A21 & Ferio XQM & \(Ebc\to(QIab\to MOac)\)\\
A22 & Complementary QE-to-QA & \(QEab\to QAab\)\\
A23 & Complementary QI-to-QO & \(QIab\to QOab\)\\
A24 & Complementary QO-to-QI & \(QOab\to QIab\)\\
A25 & Ordinary QI conversion & \(QIab\to QIba\)\\
A26 & QE subordination & \(QEab\to MEab\)\\
A27 & QI subordination & \(QIab\to MIab\)\\
A28 & QO subordination & \(QOab\to MOab\)\\
\bottomrule
\end{longtable}

There are no rejection axioms or rejection rules in this contingency calculus.
Nonacceptance is not obtained simply by reusing the L-X-M rejection apparatus
under \(Q\). The native modal immediate-inference proof above remains
available using ordinary \(Iaa\) directly instead of obtaining it from
\(LIaa\). It does not thereby prove \(LIaa\) in the new system.

From A17 with its first two subject variables identified, A1 gives
\(QAab\to QEab\). A22 gives the reverse. A18 and \(Iaa\) give
\(QAab\to QIab\); A23--A25 then supply partial complementary and ordinary
conversions. These are actual local dependencies for the collapse result,
not assumptions justified only by the phrase ``contingency conversion.''

\begin{hbtheorem}{thm:p13-johnson-q-collapse}{proved=annex:p13-johnson-q-collapse}{The two unbounded-consequent failures of Q-L-X-M}
For every well-formed formula \(r\), McCall's displayed Q-L-X-M assertion
basis proves both
\[
 QAbc\to(LAab\to r),\qquad LAbc\to(QAab\to r).
\]
These are conditional collapse results, not a claim that every formula
is an unconditional theorem.
\end{hbtheorem}

Johnson's theorems~7.1--7.2 expose this failure through actual contradictory
modal consequences. Reversal then yields acceptances missing from McCall's
comparison table, including Baroco QXM and Bocardo XQM. Agreement on a
table of selected moods therefore cannot settle the extension's behavior on
arbitrary consequents or disconnected term vocabularies
\parencite[pp.~286--288]{Johnson2004ModalSyllogisms}.

\section{Johnson's QLXM-prime repair and its exact semantics}
\label{sec:p13-johnson-qlxm-prime}

QLXM-prime deletes A20, A21, and A28 while retaining the other assertion
axioms and operations. The deletion of A20/A21 blocks the two displayed
collapse routes. A28 is deleted for a separate semantic reason: a contingent
particular is symmetric in the proposed semantics, while actual membership
in a necessary class is not. Demanding QO-to-MO subordination forces
unwanted validities such as Disamis QLQ or requires implausible constraints
on contingent sleeping and necessary humanity. This is a chosen repair,
not a proof that this is the uniquely correct historical system
\parencite[pp.~288--295]{Johnson2004ModalSyllogisms}.

\begin{hbdefinition}{def:p13-johnson-qmodel}{Nonempty-term Q-models with explicit normalization}
Use the four partition regions and actual positive abbreviation of the
J-model definition, but do \emph{not} require \(n^{+}(x)\ne\varnothing\).
Put \(q(x)=q^{+}(x)\cup q^{-}(x)\). Require \(+(x)\ne\varnothing\)
for every term, the ordinary nonemptiness needed by the restored axiom Iaa.
Require the following source BQ2--BQ8 conditions for all term parameters:
\[
\begin{array}{ll}
\mathrm{BQ2}&+(x)\subseteq n^{-}(y)\ \Longrightarrow\ +(y)\subseteq n^{-}(x),\\
\mathrm{BQ3}&+(y)\subseteq n^{+}(z),\ +(x)\mathbin{o}+(y)
                  \ \Longrightarrow\ n^{+}(x)\mathbin{o}n^{+}(z),\\
\mathrm{BQ4}&+(y)\subseteq n^{-}(z),\ +(x)\mathbin{o}+(y)
                  \ \Longrightarrow\ n^{+}(x)\mathbin{o}n^{-}(z),\\
\mathrm{BQ5}&+(z)\subseteq n^{+}(y),\ n^{+}(x)\mathbin{o}n^{-}(y)
                  \ \Longrightarrow\ n^{+}(x)\mathbin{o}n^{-}(z),\\
\mathrm{BQ6}&+(y)\subseteq q(z),\ +(x)\subseteq q(y)
                  \ \Longrightarrow\ +(x)\subseteq q(z),\\
\mathrm{BQ7}&+(y)\subseteq q(z),\quad
  (+(x)\mathbin{o}q(y)\text{ or }q(x)\mathbin{o}+(y))\\
 &\hspace{1em}\Longrightarrow\ +(x)\mathbin{o}q(z)\text{ or }q(z)\mathbin{o}+(x),\\
\mathrm{BQ8}&+(y)\subseteq +(z),\ +(x)\subseteq q(y)
                  \ \Longrightarrow\ n^{+}(x)\cap n^{-}(z)=\varnothing.
\end{array}
\]
Use the earlier six assertoric/necessary truth conditions and ordinary Boolean
connectives. Add \(QAxy,QExy\) true exactly when \(+(x)\subseteq q(y)\);
\(QIxy,QOxy\) true exactly when \(+(x)\mathbin{o}q(y)\) or
\(q(x)\mathbin{o}+(y)\). For either L or Q, normalize any even number
of categorical negations away and any odd number to one before these
clauses are used. This implements the inherited categorical DN operation.
\end{hbdefinition}

Two formal qualifications are visible. The source's printed BQ list omits
ordinary positive nonemptiness, although its soundness proof immediately uses
\(Iaa\). The present definition makes that required assumption explicit.
Also the earlier valuation definition names double-negation normalization
only for L; the Q extension needs its inherited DN treatment made explicit
under Q as well. Neither qualification restores the rejected essential-positive
existence axiom \(LIaa\).

The printed BQ7 consequent expresses the same overlap twice in reversed
order. It is retained here as the requirement \(+(x)\mathbin{o}q(z)\),
which suffices for the first disjunct of QI's truth condition. Replacing it
by the weaker, fully symmetric QI condition would change the model class;
that replacement is not silently adopted as a source correction.

\begin{hbtheorem}{thm:p13-johnson-q-nonempty}{proved=annex:p13-johnson-q-nonempty}{Why the ordinary and necessary existence assumptions differ}
The printed partition/BQ2--BQ8 conditions alone admit a structure in which
\(Iaa\) is false. The explicit ordinary-nonempty completion in
\cref{def:p13-johnson-qmodel} excludes that failure, but still admits a
model in which \(Iaa\) is true and \(LIaa\) is false for every term.
\end{hbtheorem}

\begin{hbtheorem}{thm:p13-johnson-q-soundness}{proved=annex:p13-johnson-q-soundness}{Soundness of the explicitly completed QLXM-prime basis}
With \cref{def:p13-johnson-qmodel}, every retained assertion axiom and every
assertion transformation of QLXM-prime preserves validity. Adding Johnson's
A29, \(QAbc\to(MIab\to QIac)\), also preserves validity.
This is soundness for the stated corrected model definition and restricted
language, not a completeness theorem or a proof of historical identification.
\end{hbtheorem}

A29 is not a repetition of Aristotle's accepted two-premise sample: Johnson
adds it because his semantics validates an inference with a possibility
antecedent. Its proof must split the possible-particular case using the
partition and BQ7, rather than treating MI as ordinary I. This difference is
the mechanism of the additional QMQ and MQQ constructions.

\subsection{Original Q-countermodels and the deleted subordination}
\label{sec:p13-johnson-q-countercases}

These are Johnson's original tables~16--18, now with every region explicit.
The first two source labels are \(\mathcal M_5,\mathcal M_6\); the third
is \(\mathcal M_7\). Unmentioned terms again duplicate the a-row.

\begin{longtable}{@{}llcccc@{}}
\caption{Johnson's Q-countermodels (source tables~16--18) with every region explicit}\\
\toprule Model & Term & \(n^{+}\) & \(q^{+}\) & \(n^{-}\) & \(q^{-}\)\\
\midrule\endfirsthead
\caption[]{Johnson's Q-countermodels (source tables~16--18) with every region explicit (continued)}\\
\toprule Model & Term & \(n^{+}\) & \(q^{+}\) & \(n^{-}\) & \(q^{-}\)\\
\midrule\endhead
\midrule\multicolumn{6}{r}{\emph{continued on next page}}\\
\endfoot
\endlastfoot
\(\mathcal M_5\)&a&\(\{1\}\)&\(\{2\}\)&\(\{3\}\)&\(\varnothing\)\\
&b&\(\{3\}\)&\(\{1\}\)&\(\varnothing\)&\(\{2\}\)\\
&c&\(\{1,3\}\)&\(\varnothing\)&\(\{2\}\)&\(\varnothing\)\\
\(\mathcal M_6\)&a&\(\{1\}\)&\(\{2\}\)&\(\{3,4\}\)&\(\varnothing\)\\
&b&\(\{3\}\)&\(\{1\}\)&\(\{4\}\)&\(\{2\}\)\\
&c&\(\{4\}\)&\(\varnothing\)&\(\{1,2,3\}\)&\(\varnothing\)\\
\(\mathcal M_7\)&a&\(\{1\}\)&\(\varnothing\)&\(\{3\}\)&\(\{2\}\)\\
&b&\(\{2\}\)&\(\{3\}\)&\(\varnothing\)&\(\{1\}\)\\
&c&\(\{1\}\)&\(\varnothing\)&\(\{3\}\)&\(\{2\}\)\\
\bottomrule
\end{longtable}

\begin{hbtheorem}{thm:p13-johnson-q-countercases}{proved=annex:p13-johnson-q-countercases}{Three completed native Q-models and five countercases}
The displayed structures satisfy the completed Q-model conditions.
\(\mathcal M_5\) invalidates Barbara LQQ and QO-to-MO subordination;
\(\mathcal M_6\) invalidates Felapton LQQ; and
\(\mathcal M_7\) invalidates an arbitrary-consequent instance of the first
collapse and the ordinary universal contingent conversion.
\end{hbtheorem}

In the published table~16, c has \(q^{+}(c)=\varnothing\) and
\(n^{-}(c)=\{2\}\), not the reverse. Column positions matter.
Its accompanying Barbara-LQQ paragraph prints LAac as the consequent in one
display, while computing QAac and naming LQQ. The reader uses the actual
QAac calculation. Ordinary QO-to-MO subordination fails in this model even
though complementary QI/QO equivalence holds. That is the exact semantic
reason A28 cannot be retained, rather than a generic dislike of contingency.

\subsection{Generating the additional native contingent moods}
\label{sec:p13-johnson-q-generation}

\begin{hbtheorem}{thm:p13-johnson-q-generated}{proved=annex:p13-johnson-q-generated}{Thirteen A29-based native assertions}
The amended QLXM-prime basis asserts Darii, Ferio, Datisi, Darapti, Ferison,
Felapton, Barbari, and Celaront in QMQ; and Disamis, Darapti, Dimaris,
Bramantip, and Camenop in MQQ. Their formulas use the same subject-first
categorical term placements as the corresponding ordinary mood names.
\end{hbtheorem}

The annex supplies the thirteen chains with the actual necessary local
conversions and subalternations. QMM/MQM consequences follow only where
the particular QI conclusion can be weakened to MI; QO-to-MO is no longer
available. QLL/LQL patterns come from appropriate reversals, not from
putting L in any conclusion. These restrictions explain why adding one
semantically required possibility-premise axiom does not authorize every
mixed contingent/possible combination
\parencite[pp.~299--302]{Johnson2004ModalSyllogisms}.

\section{Comparison counts and the remaining general problem}
\label{sec:p13-johnson-mccall-limits}

The later denominators differ from the earlier forty-two-cell translation
comparison. Johnson corrects McCall's contingency comparison from twenty-four
to twenty-nine disagreements among 154 first-, second-, and third-figure
cases, giving \(125/154\), about \(81\%\), agreement for Q-L-X-M.
For QLXM-prime he reports thirteen cases valid in the system but not in
Aristotle's classification and nine in the reverse direction: \(132/154\),
about \(86\%\), agreement. Seventeen discrepancies concern reversal pairs.
In a classical assertion system one cannot accept one conditional and reject
its tautologically equivalent reversal just to match an inconsistent historical
table. The other five discrepancies concern Darapti XQQ/LQQ, Festino QXM,
Celarent QLX, and Felapton XQM. Their historical evaluation and the deletion
of Ferio XQM remain part of Johnson's interpretation, not an empirical
performance score \parencite[pp.~288, 303]{Johnson2004ModalSyllogisms}.

For his two-premise tally, \([A]\) means an A-or-E categorical position and
\([I]\) an I-or-O position. Thus QQQ \([A][A][A]\) expands to eight
quality combinations, whereas QXQ \([A]I[I]\) expands to four. These brackets
are enumeration notation, not new categorical operators. The source's grouped
counts are:

\begin{tabularx}{\textwidth}{@{}X r X r@{}}
\toprule Group & Count & Group & Count\\\midrule
Noncontingent apodeictic cases &333& QQQ &64\\
QXQ and QLQ together &40& XQQ and LQQ together &32\\
QXM &34& XQM &20\\
QLX &24& LQX &30\\
QLM &46& LQM &46\\
QMQ &16& MQQ &16\\
QMM &8& MQM &8\\
QLL &8& LQL &8\\
QQM &48&&\\
\bottomrule
\end{tabularx}

Their sum is 781 cases in thirty-five general mode patterns, as reported in
the original pp.~303--305 tally. This source count does not establish
completeness for arbitrary premise sets, nor does it independently prove every
entry in the source's tables~13,15,23. The complete arguments supplied here
have precisely the scopes of their stated results.

Johnson ends with the proposed extension to a Smiley-type procedure for
contingent syllogisms with more than two premises satisfying the chain
condition. He proposes inconsistent cycle patterns and a negated-Q prefix;
he does not present that general procedure as already established in this
chapter. He also notes that deleting Barbara XQM would simplify a still
weaker system. The closing alternatives therefore remain alternatives, not
theorems manufactured from the 781-case tally
\parencite[pp.~305--306]{Johnson2004ModalSyllogisms}.

\begin{proofboundary}{def:p13-johnson-mccall-proof-boundary}{Scope of the displayed calculi and proofs}
The rejection deductions use the displayed partial R1--R4 basis and its
declared transformations. They do not establish rejection completeness.
The term-region results prove the stated countercases and soundness for
the explicitly completed QLXM-prime model definition; they do not prove
general L-X-M or QLXM-prime completeness. The AP/AS chain procedures and
the source's table classifications are separate results, not consequences
of the finite collection of derivations and countermodels given above.
In particular, the 781-case tally is not a decision theorem for arbitrary
premise sets or for contingent chains with more than two premises.
\end{proofboundary}
